% 第 42 章　经典物理遇到的几场灾难
\chapter{经典物理遇到的几场灾难}

一九〇〇年春天，开尔文勋爵在一次著名的演讲里说：物理学的大厦已经基本落成，剩下的只是些修饰工作，天边只飘着两朵小乌云。他说这话并非狂妄——牛顿力学解释了从苹果到行星的运动，麦克斯韦方程统一了电、磁和光，热力学和统计物理正在解开热之谜。两座奖杯擦得锃亮，谁好意思说大厦要塌？

结果，那两朵小乌云里各藏着一场风暴。一朵吹出了相对论，那是第 38 章到第 40 章的故事。另一朵——关于热的东西为什么会发光——吹倒了更多人，吹出了量子力学。这一章，我们来清点这场风暴的登陆过程：经典物理在五年间接连输掉了五场仗，每一场都输在实验事实上，无处可躲。

\step{反常识开场}

冬天晚上，看一眼厨房里的电陶炉。刚通电时，炉盘摸上去已经烫了（别真摸），可它不发光，看上去还是黑的。几分钟后，它泛起暗红色；火力调大，变成橙红；如果你有一只能达到更高温度的工业炉，会看到它继续变成亮黄、甚至白得刺眼。

铁匠早就把这件事用成了手艺：不看温度计，只瞟一眼炉子里铁块的颜色，就知道温度够不够锻打。请注意这句话里藏着的反常之处——铁匠看的是铁块，可如果他往炉子里同时扔进去一块陶瓷、一截炭、一块耐火砖，烧到同样的温度，它们发出的光，颜色几乎一样。

颜色只认温度，不认材料。

这不是工匠的错觉，而是可以在实验室里反复验证的事实：任何烧到一千开的物体都泛红光，烧到一千五百开都转橙黄，太阳表面五千八百开，正好白亮。物体的成分、硬度、产地统统退场，只剩下“温度”一个变量在指挥颜色。

现在轮到理论出场了。一九〇〇年，物理学家手里握着当时最好的两件武器——麦克斯韦的电磁理论和统计物理——来算这道题：一个热物体应该在各种颜色的光上各发多少？算出来的答案工整、漂亮，而且错得离谱。那个理论预言：任何热物体发出的辐射，频率越高能量越大，一路冲向紫外线、X 射线、伽马射线，总能量是无穷大。换句话说，按当时最好的物理，你每次站在烤箱旁边，都应该被致命的短波辐射当场击倒。

你显然没有被击倒。烤箱每天都开，没有人因此进医院。理论错得如此彻底，错在什么假设上？这一章就要把这个假设揪出来。揪出来以后你会发现，它顺手还解释了另外两件怪事：为什么荧光灯摸上去温温的，却能发出亮白光；以及——你身上此刻正不断向外发射一种看不见的光。

\step{先猜一猜}

照例，拿一张纸，先写下你的直觉答案。本章结尾请回来给自己打分。

第一题：一块铁和一块陶瓷放进同一个炉子，烧到同样的一千度。它们发光的颜色，是铁更红、陶瓷更红，还是基本相同？

第二题：实验室里有一块金属板。用一束很强的红光长时间照射它，和用一束很微弱的紫光短时间照射它，哪一种做法更可能把金属里的电子“打”出来？或者，你认为两者都行，只是快慢不同？

第三题：把炉子的温度一升再升，物体发光的颜色会一路红、橙、黄、白变化下去。这条路有没有尽头？温度足够高时，光会变成蓝色吗？蓝色再往上呢？

写好了吗？如果你对第二题的答案是“红光够强够久总能行”，你和一九〇五年之前所有的物理学家想的一样——而这恰恰是错得最脆生的一题。

\step{动手实验}

\begin{experiment}[用一张废 CD 拆开两种白光]
“光里面到底有哪些颜色”听起来需要昂贵的仪器，其实一张报废光盘就够。

\textbf{做法：}找一张废旧 CD 或 DVD（盘面无标签、能反光的）。在暗室里点亮一盏白炽灯或调光台灯，让灯光斜照在光盘有彩虹纹的一面，眼睛凑近盘面，从不同角度寻找反射光。你会在某个角度看到一条明亮的彩色光带。记下它。然后把光源换成节能灯（荧光灯）或手机纯白屏幕，重复同样的动作。

\textbf{你会看到：}白炽灯的光带是一条连续的彩虹，红橙黄绿蓝紫一个不缺，彼此平滑过渡。节能灯的光带却断成了几截：只有几条孤零零的亮带，其中绿色那条格外扎眼，带与带之间是大片空缺。手机白屏的光带则是另一番模样：一团窄窄的蓝，加一团宽宽的黄，中间也缺着什么。

\textbf{为什么：}光盘上每毫米刻着六百多圈细槽，槽间距约 \(1.6\,\mu\mathrm{m}\)，和光的波长同量级，整张盘就是一面巨大的衍射光栅。不同颜色的光被这些刻槽“弹”向不同的角度（衍射的原理见第 35 章），于是一束白光里藏着哪些颜色，就被摊开成了一张清单。白炽灯是“烧热了发光”，清单连续；节能灯是气体放电发光，清单上只有几行。

\textbf{想一想：}这个实验没有解释任何事，它只是把两件事并排摆在了桌上：为什么“热的”光源交出连续清单，“不热的”光源交出稀疏清单？请记住你亲眼看到的这几条亮带，它们后面会变成原子寄给你的密码信。
\end{experiment}

\step{旧模型遇到麻烦}

现在，把镜头推回一九〇〇年前后，看经典物理是怎样一场接一场输掉的。

\textbf{灾难一：热辐射的总账算成了无穷大。}

先补一个事实：任何温度高于绝对零度的物体都在发光，这叫热辐射。温度低时，光的主体在红外波段，你的眼睛看不见，所以刚通电的炉盘“黑着”；但蛇的信子、热成像仪看得见——你此刻就在向房间各处发射红外线，这一点后面还要回来细说。

经典理论算这笔账的思路完全正当。把炉腔里的电磁场想象成一大堆“振子”：每种频率的光波对应一种驻波模式，就像一间屋子里同时挂着无数根不同调子的琴弦。统计物理（第 31 章）有一条久经考验的规则——能量均分：达到热平衡时，每种模式平均分到同样大的一份能量 \(k_{\mathrm{B}}T\)，\(k_{\mathrm{B}}\) 是玻尔兹曼常数，\(T\) 是温度。这条规则在气体分子身上验证过无数次，没理由怀疑它。

然后灾难来了。波长可以任意短，频率可以任意高，高频模式有无穷多个，而每个模式都要领走一份 \(k_{\mathrm{B}}T\)。无穷多个非零的数相加，结果是无穷大。沿着这条思路写出的公式（瑞利—金斯公式）在长波端和实验吻合得很好，可频率一高就发疯：它预言辐射强度随频率一路飙升，永不回头。后人给它起了个名字——紫外灾难。

实验测到的真实曲线是什么样子？从低频升起，在某个频率达到一个峰，然后乖乖落回零。太阳表面约 \(5800\,\mathrm{K}\)，峰值在约 \(500\,\mathrm{nm}\)，正是可见光；一千开的炉丝峰值在红外，只有红色的尾巴尖漏进可见区，所以它暗红；你的体温 \(310\,\mathrm{K}\)，峰值在约 \(10\,\mu\mathrm{m}\) 的远红外。理论在长波端全对，在短波端全错。错得那么干脆，说明出事的不是某个系数，而是某个根本假设。

\textbf{灾难二：光打电子，打得完全不讲道理。}

一八八七年，赫兹在验证电磁波存在的实验里顺手注意到一个怪现象：用紫外光照射金属电极，火花更容易跳出来。十几年后，勒纳德等人把它做成了精密实验：光照射金属表面，确实能把电子打出来，这叫光电效应。

如果光是一种波——这是当时铁板钉钉的结论——那么下面三个预言看起来不容置疑：

\begin{itemize}[leftmargin=1.8em]
\item 光越强，波的能量越大，打出的电子应该跑得越快；
\item 任何频率的光，只要照得够强、够久，能量攒够了总能打出电子；
\item 光很弱时，电子要等能量慢慢攒够才出来，应该有明显的延迟。
\end{itemize}

实验的回答是三连否。第一，电子的最大能量只取决于光的颜色：紫光打出的电子比蓝光的快，蓝光的比绿光的快；把红光增强一万倍，打出的电子一个也不多、一个也不快。第二，对每种金属都存在一个截止频率：光的颜色比它“红”，哪怕用探照灯照上一年，一个电子也出不来；颜色比它“紫”，微弱到几乎看不见的光也立刻打出电子。第三，只要颜色够“紫”，电子出来的速度快得测不出延迟——后来确认响应时间小于十亿分之一秒量级。

这个反例值得停下来体会，因为它直接扇在直觉上。一束很强的红光，输给了一束很弱的紫光。“量大”败给了“频率高”。如果用经典图像估算延迟，会更难看：普通强度的光照在金属上，按波的图像，一个原子尺度的区域接收能量的功率大约是 \(10^{-20}\,\mathrm{W}\) 量级，而打出一个电子需要约几电子伏特、即 \(10^{-19}\,\mathrm{J}\) 量级的能量，照理要攒上几分钟乃至几小时。实验说：不到十亿分之一秒。数量级差了十亿倍。

\textbf{灾难三：原子根本不应该存在。}

第三场灾难更彻底。氢气装进放电管，通电，发光，用棱镜或光栅拆开——不是彩虹，而是几条特定颜色的亮线。早在一八八五年，中学数学教师巴耳末就发现氢的这几条谱线波长满足一个简单得可疑的经验公式。可没人知道为什么：原子为什么会发光？为什么只发这几种颜色？

更糟的是稳定性问题。卢瑟福的散射实验（详见第 51 章）表明原子有一个极小极重的核，电子在核外。按当时唯一可能的图像，电子绕核运动——这是加速运动，而麦克斯韦理论（第 24 章）说得很清楚：加速的带电粒子必然辐射电磁波，不断损失能量。算下来，电子应该在约 \(10^{-10}\) 秒内沿螺旋线掉进原子核。也就是说，经典物理预言：原子会在百亿分之一秒内坍缩，物质不存在，你也不存在。

三场灾难指向同一个方向：经典理论不是哪一处算错了，而是它对“能量可以连续地、任意小份地流动”这件事的默认，在微观世界里不成立。

\step{发明新概念}

历史上有趣的一幕是：修补工作是从一个“作弊”开始的。

\textbf{普朗克的补丁（一九〇〇年）。}为了凑出那条先升后降的辐射曲线，普朗克试了一个他自己都觉得别扭的假设：频率为 \(f\) 的振子和外界交换能量时，不能想交换多少就交换多少，只能一整份一整份地交换，每份的大小是
\[
E = hf,
\]
其中 \(h = 6.63\times 10^{-34}\,\mathrm{J\cdot s}\) 是一个全新的自然常数，今天叫普朗克常数。能量像硬币，有最小面额，面额和频率成正比。

这个补丁为什么能消灭紫外灾难？高频振子的一份能量 \(hf\) 太大，而热运动平均只能拿出约 \(k_{\mathrm{B}}T\) 的零花钱。频率高到一定程度，一份都凑不齐，于是高频振子根本领不到能量——它们被“冻结”了。无穷多个高频模式，绝大多数颗粒无收，总账立刻变得有限。曲线先升后降，和实验分毫不差。

这个比喻值得说完整。它像自动售货机：只收整枚硬币，标价太高的商品，零钱再多也买不走。它不像关小的水龙头：不是能量流得慢了，而是交易的最小单位本身被抬高了。普朗克本人很多年都不喜欢这个补丁，觉得它只是数学技巧。真正把它当真的是一个二十六岁的专利局职员。

\textbf{爱因斯坦的光子（一九〇五年）。}爱因斯坦说：别只怪振子，干脆承认光本身就是这样的。频率为 \(f\) 的光，在与物质交换能量时，永远以 \(hf\) 为一整包出手，这一包就叫光子。光电效应的三连否立刻变成了三连通：

一个电子一次只吞一个光子。光子先替电子付一张“离开金属的门票”——逸出功 \(W\)，每种金属门票价格不同——剩下的钱变成电子的动能：
\[
hf = W + E_{\mathrm{k}}.
\]
逐条对账。颜色不够“紫”（\(hf < W\)）：一张门票都买不起，光再强，也只是亿万个穷光子排长队，而电子一次只接待一个——所以有截止频率，强度救不了场。颜色够：光子一到就成交，没有攒钱的过程——所以没有延迟。光变强：来的光子多了，成交的电子多了，但每个电子付完门票剩下的钱不变——所以强度只决定电子数量，不决定电子速度。每一条都严丝合缝。

\textbf{玻尔的能级（一九一三年）。}面对离散光谱和“原子为什么不坍缩”，玻尔的办法同样不讲理：直接规定。原子内部只允许存在某些特定的能量状态，叫定态；处于定态的原子不辐射，麦克斯韦那套“加速必辐射”在这里被叫停；辐射只发生在原子从高能级跳到
低能级的瞬间，发出的光子能量恰好等于两级之差：
\[
hf = E_{\text{高}} - E_{\text{低}}.
\]
能级是离散的，能级差是离散的，光子能量自然是离散的——光谱当然只剩几条亮线。巴耳末那几条神秘的谱线波长，代入氢的能级 \(E_n = -13.6\,\mathrm{eV}/n^2\)（\(n = 1, 2, 3, \dots\)）后一条不差地复现。这里必须立刻立一块护栏：玻尔模型里的电子不是沿确定路线绕核飞的行星，定态不是一条可以画出来的轨道，“电子在某圈轨道上跑”是旧词残留，别把它当图像；这个模型为什么部分正确、边界在哪里，要到第 44 章和第 52 章才有真正的答案。

\textbf{光子的动量：康普顿散射（一九二三年）。}如果光真的是一包一包的东西，那它除了能量还该有动量。第 40 章讲过相对论的能量—动量关系：静质量为零的粒子满足 \(p = E/c\)，于是光子的动量是
\[
p = \frac{h}{\lambda}.
\]
康普顿用 X 射线轰击石墨里的电子。如果 X 射线只是波，电子被波推着晃，散射出来的波频率不该变。实验却测到：散射后的 X 射线波长变长了，且变长量随散射角变化，精确等于“一个光子撞上一个静止电子、按能量动量守恒弹开”算出的结果——偏转九十度时波长约变长 \(0.0024\,\mathrm{nm}\)，头对头反弹时变长两倍。光和电子的这一撞，是按两粒台球碰撞的账本平的。光子有动量，不再是比喻，而是测量值。

\textbf{电子的波长：德布罗意与电子衍射（一九二四至一九二七年）。}最后一场颠覆方向相反。德布罗意问：波能被数成一份份的粒子，那一向被当成粒子的东西，会不会其实是波？他猜任何动量为 \(p\) 的粒子都伴随着一个波长
\[
\lambda = \frac{h}{p}.
\]
三年后，戴维孙和革末把电子束射向镍晶体，G. P. 汤姆孙让电子穿过薄金属膜，都看到了明暗相间的衍射图样——那是波的签名，粒子无论如何画不出来。电子，这个在阴极射线管里被电场推着走、被当成小颗粒三十年的东西，会衍射。

到这里，五场灾难的战场都清点完了。请不要把战果总结成“光和电子有时装成波、有时装成粒子”——那种说法是把两个借来的词当成两种 costume。正确的说法是：从宏观世界借来的“波”和“粒子”两个词，单独哪一个都装不下实验事实；光子和电子从来只有一种存在方式，只是我们的旧词汇表不够用。够用的那个词汇表长什么样，是第 43 章的主题。

\step{一句话模型}

\begin{oneline}
光和物质交换能量时从不零存零取，永远以 \(E = hf\) 的整包成交；经典物理的所有灾难，都来自“能量可以无限细分”这个从未被检验过的默认。
\end{oneline}

\step{数学翻译器}

直觉已经就位，现在把它翻译成可以计算的式子，并且每条式子都当场用极端情形检查一遍。

\textbf{第一式：\(E = hf\)。}它描述：频率为 \(f\) 的光，每包能量是多少。为什么这样写：正比是最简单的关系，而黑体辐射曲线逼出来的正是它。检查特殊情况——频率为零：\(E = 0\)，不振就没有包可发，合理。频率极大：伽马射线 \(f \sim 10^{20}\,\mathrm{Hz}\)，单包能量约 \(7\times10^{-14}\,\mathrm{J}\)，顶得上约二十万个可见光光子挤在一起，所以伽马射线防护要用铅板而不是玻璃。频率极小：调频广播 \(f \sim 10^{8}\,\mathrm{Hz}\)，单包能量约 \(7\times10^{-26}\,\mathrm{J}\)，小到任何接收机都数不出来——这就是为什么收音机里的电磁波看上去完全是连续的波。量子性不是不出现，是包太小，被平均掉了。

\textbf{第二式：光电方程 \(hf = W + E_{\mathrm{k}}\)。}它描述：一个光子把能量交给一个电子后，钱怎么分。检查特殊情况——\(hf\) 恰好等于 \(W\)：电子动能 \(E_{\mathrm{k}} = 0\)，恰好在门槛上，对应的频率 \(f_0 = W/h\) 就是截止频率。\(hf < W\)：动能算出来是负数，物理上不可达，一个电子也出不来，公式自己宣告了自己的边界。光强翻倍：式子里根本没有“强度”这一项，单个电子的动能纹丝不动，翻倍的只是每秒成交的笔数。

\textbf{第三式：巴耳末公式。}氢的可见光谱线满足
\[
\frac{1}{\lambda} = R\left(\frac{1}{2^2} - \frac{1}{n^2}\right), \qquad n = 3, 4, 5, \dots
\]
其中 \(R \approx 1.097\times10^{7}\,\mathrm{m^{-1}}\)。检查特殊情况——\(n = 3\)：\(\lambda \approx 656\,\mathrm{nm}\)，正是氢那道著名的红；\(n \to \infty\)：\(\lambda \to 365\,\mathrm{nm}\)，谱线越挤越密，收敛到一个极限波长，这个“系列限”对应电子刚好挣脱原子的能量。一条纯经验公式，被能级差 \(hf = E_{\text{高}} - E_{\text{低}}\) 完整解释——这不是巧合，是能级存在的收据。

\textbf{第四式：德布罗意波长 \(\lambda = h/p\)。}先看人。六十公斤的人以每秒一米走路，\(\lambda = h/(mv) \approx 10^{-35}\,\mathrm{m}\)，比原子核还小二十个数量级。所以你穿过门框时不会发生衍射——不是原理不允许，是波长不够格。再看电子：被一百伏电压加速的电子，动量约 \(5.4\times10^{-24}\,\mathrm{kg\cdot m/s}\)，波长 \(\lambda \approx 0.12\,\mathrm{nm}\)，恰好和晶体里的原子间距同量级——晶体对电子束来说，就是一张天然光栅。公式两边一对照，电子衍射实验为什么用晶体，一目了然。

\textbf{一个带真实数据的估算：一盏灯每秒发多少个光子？}取一只 \(100\,\mathrm{W}\) 的白炽灯，它只有约百分之五的电能变成可见光（其余变成热，这正是白炽灯被淘汰的原因），即约 \(5\,\mathrm{J/s}\)。可见光取平均波长 \(550\,\mathrm{nm}\)，单光子能量
\[
E = hf = \frac{hc}{\lambda} \approx \frac{6.63\times10^{-34}\times3.0\times10^{8}}{5.5\times10^{-7}} \approx 3.6\times10^{-19}\,\mathrm{J}.
\]
每秒发出的光子数约为 \(5 / 3.6\times10^{-19} \approx 1.4\times10^{19}\) 个。一万万亿个光子每秒钟从你台灯里涌出来。现在你知道“连续的光”是什么了：不是不分包，是包太多、太小，任何仪器和眼睛都只能看到它们的洪流，看不到单个的雨点。

\begin{mathbridge}
\textbf{紫外灾难的定量形状，以及普朗克曲线怎样收场。}

瑞利—金斯的经典结果：腔内频率在 \(f\) 附近的辐射能量密度为
\[
u(f) = \frac{8\pi f^2}{c^3}\,k_{\mathrm{B}}T.
\]
因子 \(f^2\) 来自“频率越高、驻波模式越多”，\(k_{\mathrm{B}}T\) 来自能量均分。把它对所有频率积分，\(\int_0^\infty u(f)\,df\) 发散——这就是紫外灾难的原文，灾难不是修辞，是积分号下的无穷大。

普朗克把每个振子的平均能量从 \(k_{\mathrm{B}}T\) 换成
\[
\bar{E} = \frac{hf}{e^{hf/k_{\mathrm{B}}T} - 1},
\]
得到普朗克分布。检查两个极限，看新公式如何两头讨好——低频端 \(hf \ll k_{\mathrm{B}}T\)：把指数展开，\(e^{x} - 1 \approx x\)，于是 \(\bar{E} \approx k_{\mathrm{B}}T\)，自动退回经典的均分结果；高频端 \(hf \gg k_{\mathrm{B}}T\)：分母里的指数爆炸式增长，\(\bar{E}\) 被指数压向零，高频振子集体冻结，积分收敛。经典理论没有被推翻，它被证明是量子理论在 \(hf \ll k_{\mathrm{B}}T\) 时的近似——这正是它骗过了两百年的原因。

对普朗克分布求峰值位置，得到维恩位移定律：\(\lambda_{\max} T \approx 2.9\times10^{-3}\,\mathrm{m\cdot K}\)。温度翻倍，峰值波长减半。开场的颜色序列——暗红、橙红、黄、白——全在这一个式子里。
\end{mathbridge}

\begin{challenge}
\textbf{康普顿散射：用守恒律算出波长的改变。}

把散射当成一次弹性碰撞：光子带着能量 \(hc/\lambda\) 和动量 \(h/\lambda\) 撞上一个静止的电子（质量 \(m_{\mathrm{e}}\)），碰后光子沿角 \(\theta\) 飞出、波长变为 \(\lambda'\)，电子带着反冲能量和动量离开。写下两条守恒——能量守恒，以及动量在两个方向上的分量守恒（动量是矢量，要分方向写）。电子的能量用相对论关系 \(E^2 = p^2c^2 + m_{\mathrm{e}}^2c^4\)。消去电子的能量和动量，剩下的就是
\[
\lambda' - \lambda = \frac{h}{m_{\mathrm{e}}c}(1 - \cos\theta).
\]
检查特殊情况——\(\theta = 0\)：\(\cos\theta = 1\)，波长不变，没撞就没变，合理；\(\theta = 180^\circ\)：改变量取最大值 \(2h/(m_{\mathrm{e}}c) \approx 4.9\times10^{-12}\,\mathrm{m}\)，头对头反弹时光子损失能量最多。常数 \(h/(m_{\mathrm{e}}c) \approx 2.4\times10^{-12}\,\mathrm{m}\) 叫电子的康普顿波长。注意这个推导从头到尾只用了“光子有能量 \(hf\)、动量 \(h/\lambda\)”加守恒律——它既是公式的来源，也是光子动量最直接的证据。
\end{challenge}

\step{模型的边界}

新模型威力巨大，但每一条都有写在包装上的适用范围。

第一，光子不是一颗微缩弹珠。“粒子性”只在一处兑现：能量以 \(hf\) 整包交换。光子没有一条可以画出来的小球轨迹，它的传播、干涉、绕障碍，仍然由波动（更深处是量子场论）描述。把它想象成硬球，双缝实验会立刻让你出丑。

第二，波粒二象性不是变身术。光子和电子从不在“波的形态”和“粒子的形态”之间切换；是“波”和“粒子”这两个经典类别各自都不完整。你设计什么样的实验，就看到哪一类现象浮出水面，但被测的对象自始至终只有一个。第 43 章会用逐个粒子的双缝实验把这句话钉死。

第三，\(E = hf\) 里的 \(f\) 是光作为一种周期过程的频率，别把它读成“小球每秒抖动 \(f\) 次”。公式说的是交换的账本，不是小球内部的构造。

第四，玻尔模型是过渡品，不是终点。它精确命中氢的谱线，却在氦原子上失灵，也算不出谱线的亮度。它的价值在于第一次把“能级”请进了物理，而它自己为什么不自洽，要等薛定谔方程（第 44 章）来回答。

第五，也是最容易被误用的一条：量子不是对经典的否定，而是给经典划了地盘。凡是 \(hf \ll k_{\mathrm{B}}T\)、凡是动量与尺寸的乘积远大于 \(h\) 的场合，量子公式自动退回经典公式。篮球、行星、微波炉里的光波，照旧用经典物理，一分不差。典型误用清单：以为“光够强就什么电子都能打出来”（错，频率是门槛）；以为“看颜色能测任何光源的温度”（错，荧光灯和 LED 的颜色由能级差决定，与温度无关——下一段就用到这一条）；以为“量子效应意味着宏观世界也充满不确定性”（错，你的德布罗意波长是 \(10^{-35}\) 米）。

\step{回头解释开场现象}

现在回到厨房那只炉子，把每个悬念逐个关掉。

为什么颜色只认温度、不认材料？热辐射是物体内部电荷的热运动发出的电磁波，达到热平衡后，辐射的谱形由普朗克分布给出，而普朗克分布里唯一的宏观参数就是温度。材料的影响（发射率）只改变“发得多还是少”，不改变“峰在哪里”。铁、陶瓷、炭烧到同一温度，亮度可以有差别，颜色却基本一致——铁匠的手艺，本质是拿眼睛当一台粗糙的光谱仪，读的是同一条普朗克曲线。

为什么先红后黄再白？维恩位移定律：温度升高，峰值波长等比例缩短。一千开时峰在深红外，只有红色波段的尾巴漏进人眼，所以暗红；一千五百开，橙黄大量涌入；太阳表面的五千八百开，峰值 \(500\,\mathrm{nm}\) 正好落在可见光正中央，全部颜色一起涌来，混在一起就是白。调光台灯调暗时电流变小、灯丝变凉，整条曲线向红端滑动，灯光随之变黄变红——你家里的旋钮，就是在亲手拨动维恩位移定律。

紫外灾难为什么没发生？高频振子一份能量就要 \(hf\)，热运动的零钱买不起，它们被冻结在账外。烤箱安静，是因为普朗克常数不为零。

荧光灯为什么不烫也发白光？它根本不在“热发光”的游戏里。灯管里汞蒸气放电，电子撞击汞原子，把原子送上高能级；原子跳回低能级时按 \(hf = E_{\text{高}} - E_{\text{低}}\) 发出几条特定的谱线，其中最强的在紫外；管壁上的荧光粉吸收这些紫外光子，再按自己的能级差把它们换成可见光的宽谱。你在 CD 实验里看到的那几条分立亮带，就是汞原子能级差的直接照片。它的颜色由能级写死，摸上去只有温温的几十度——“看颜色知温度”对它彻底失效，因为那张颜色的价目表从来不是温度开的。

还有你身上那束看不见的光。体温 \(310\,\mathrm{K}\)，维恩位移给出峰值波长约 \(10\,\mu\mathrm{m}\)，在远红外。你每平方米皮肤大约以几百瓦的功率向外辐射（同时接收环境的辐射，净账小得多）。这不是修辞，而是产业：热成像仪的核心就是一片对 \(10\,\mu\mathrm{m}\) 光子敏感的探测器，用光电效应的现代后代——半导体探测器——把你发出的每包红外能量数成一个像素。消防搜救、体温筛查、夜视设备，全是本章公式的直接应用。手机摄像头里的 CMOS、屋顶的太阳能板、物理实验里的光电倍增管，也都是同一本账：光子到，电子动，一包一笔。

\step{脑洞实验与挑战题}

\begin{thought}[把普朗克常数的旋钮慢慢拧大]
量子效应在日常里隐身，唯一的原因是 \(h\) 太小。做个思想实验：想象宇宙有个旋钮，能把 \(h\) 慢慢调大，拧到多大，你会亲眼看见世界“变量子”？

拿台球来算。一只台球约 \(0.2\,\mathrm{kg}\)，以 \(1\,\mathrm{m/s}\) 滚向五厘米宽的袋口。要让它的德布罗意波长达到袋口的尺寸——\(\lambda = h/(mv)\)，反解出 \(h = mv\lambda \approx 0.2\times1\times0.05 = 10^{-2}\,\mathrm{J\cdot s}\)，是真实值的约 \(10^{32}\) 倍。拧到那一格：台球滚向袋口时会像波一样散开衍射，进不进袋变成概率问题；可见光光子一个个像子弹，一盏灯每秒只发出寥寥几颗，“连续的光”消失；百米赛跑的成绩再也测不准。

这个实验像什么：它像拧收音机的音量，帮你听清一个被压低的频道——量子规则其实一直在响，只是音量太小。它不像什么：它不像对真实世界的改造方案。真把 \(h\) 调大，原子尺寸、光谱、化学键全部重写，“台球”和“你”根本不会存在。旋钮只能想象，不能真拧。
\end{thought}

\begin{puzzles}
\item 两只完全相同的白炽灯泡，一只接足电压发白光，一只调低电压发暗红光。哪一只每秒发出的光子更多？\textbf{思路提示：}别只比较功率。红光光子单价低，同样一块钱能买更多个；先把两只灯的总功率想成可调的，再数一数“个数”和“单价”两个变量谁变得快。
\item 某金属恰好能被绿光打出电子。现在改用强度加倍的红光，再改用强度减半的紫光，分别会发生什么？\textbf{思路提示：}光电方程里有没有“强度”这一项？强度出现在物理过程的哪个环节——单个电子的能量，还是每秒成交的笔数？
\item 太阳辐射的峰值恰好落在可见光波段，人眼又恰好对这一段最敏感。这是巧合吗？\textbf{思路提示：}把因果倒过来问——眼睛是在什么样的光环境里演化出来的？哪一侧是原因，哪一侧是结果？再想一层：如果太阳表面温度变成一万开，地球生命可能“看”什么颜色？
\item 电子显微镜能看见病毒，光学显微镜不能。为什么？\textbf{思路提示：}显微镜的分辨本领被波长卡住，细节小于波长就糊成一团。病毒约 \(100\,\mathrm{nm}\)，可见光约 \(500\,\mathrm{nm}\)；再用 \(\lambda = h/p\) 估一估被一万伏加速的电子波长，和病毒比一比。
\end{puzzles}

五场灾难清点到此为止。我们手里多了一把新钥匙：能量以 \(hf\) 为包交换，粒子带着波长 \(h/p\) 行走。可这把钥匙刚刚插进锁孔——如果电子真的是“波”，那一个电子同时穿过两条缝时到底发生了什么？那不是灾难，是比灾难更奇怪的东西。第 43 章见。
